%%%%%%%%%%%%%%%%%%%%%%%%%%%%%%%%%%%%%%%%%%%%%%%%%%%%%%%%%%%%%%%%%%%%%%%%
%
% Two-page extended abstract
% From Human Videos to Dexterous Robot Capability workshop (CoRL 2026)
% Formatting: NVIDIA tech report.
%
% Limit: 2 pages of main text; references, acknowledgments and
% appendices do not count.
%
% Build:  pdflatex main && bibtex main && pdflatex main && pdflatex main
%
%%%%%%%%%%%%%%%%%%%%%%%%%%%%%%%%%%%%%%%%%%%%%%%%%%%%%%%%%%%%%%%%%%%%%%%%

\pdfoutput=1

%%% == IMPORTANT ==
%%% For anonymized submission (no authors, logo, copyright or
%%% acknowledgments), use this:
%\documentclass[10pt, logo, copyright, anonymous]{nvidiaextabstract}

%%% For camera-ready, use this:
\documentclass[10pt]{nvidiaextabstract}

%%% Load required packages here (many are loaded by the class already:
%%% amsmath, amssymb, graphicx, booktabs, xcolor, hyperref, natbib,
%%% balance, enumitem, caption).
\usepackage{microtype}
\usepackage{multirow}
\usepackage{subcaption}
\usepackage{xspace}

%%%%%%%%%%%%%%%%%%%%%%%%%%%%%%%%%%%%%%%%%%%%%%%%%%%%%%%%%%%%%%%%%%%%%%%%

%%% Venue: printed in the first-page footer. Comment out to drop it.
\venue{From Human Videos to Dexterous Robot Capability (CoRL 2026)}

%%% Logo in the top-left corner of page 1 (shown with the `logo' class
%%% option). Options in [] go to \includegraphics; default is the NVIDIA
%%% logo at width=90pt. The header is about 40pt tall, so keep height below that.
%\logo[height=28pt]{assets/mylogo.png}

%%% Printed instead of the authors when the `anonymous' option is set.
\submissionID{maniphys.ai}

\title{Per-Episode Policy Selection with an Object-Orientation Reward for Video-to-Robot Grounding}
\shorttitle{maniphys.ai: Track 2 Technical Report}

%%% Authors: to be filled in (names, order, affiliations, emails).
\author[1]{Author list to be provided}
\affil[1]{maniphysAI lab}

\begin{abstract}
\textbf{Abstract:}
We describe the maniphys.ai entry to Track 2 (Robotic Grounding) of the V2D Challenge, which covers all three input tiers with the Unitree G1 and Dex3 hands in the official FlashCHORD simulator.
We train one FlashSAC policy per episode and per tier and change the official recipe in only two places: a tighter object-tracking setting and one extra reward term on the object orientation error.
Rules fixed before each experiment decide, from a 4096-world evaluation against the tier's own reference, which candidate policy enters a submission.
For the noisy and video-only tiers we add a clean-room trajectory denoiser and a monocular human-object reconstruction chain with frame-0 alignment.
The orientation term raised 4096-world ADD-AUC on 11 of 12 tracked Tier 2 and Tier 3 episodes by 0.08 to 0.37.
On the public boards (3 November) the entry ranks first on Tier 3 SP-SR.
\end{abstract}

%%%%%%%%%%%%%%%%%%%%%%%%%%%%%%%%%%%%%%%%%%%%%%%%%%%%%%%%%%%%%%%%%%%%%%%%

%%% Author-defined commands.
\definecolor{NvidiaGreen}{RGB}{118, 185, 0}
\newcommand{\method}{\textsc{Method}\xspace}
\newcommand{\rtho}{\textsc{r3o}\xspace}

%%%%%%%%%%%%%%%%%%%%%%%%%%%%%%%%%%%%%%%%%%%%%%%%%%%%%%%%%%%%%%%%%%%%%%%%

\begin{document}

\maketitle

%%%%%%%%%%%%%%%%%%%%%%%%%%%%%%%%%%%%%%%%%%%%%%%%%%%%%%%%%%%%%%%%%%%%%%%%

\section{Introduction}
\label{sec:intro}

Track 2 asks for a robot policy that reproduces each of 30 human-object interaction episodes, such as righting a fallen bar stool or handing a basket from one hand to the other.
The input comes in three tiers: clean trajectories (Tier 1), trajectories with synthetic noise (Tier 2), and a single exocentric video (Tier 3).
The tasks are whole-body and bimanual: one policy has to walk, make contact and carry the object without an assist force at the end of training.

We keep the official FlashCHORD training recipe, built on FlashSAC, a soft actor-critic variant~\cite{haarnoja2018sac}, as the base and treat the problem as 90 independent tracking problems (30 episodes, 3 tiers).
Most of the effort went into deciding, per episode, which variant to trust: one episode moves a board by only 1/30, and single-seed results vary widely.
Our contributions are:
\begin{itemize}
  \item an object-orientation reward term that raised ADD-AUC on most tracked episodes on all three tiers without other changes;
  \item a per-episode selection procedure with swap rules fixed before each experiment, and an upload gate that discounts own-reference gains by a measured transfer ratio;
  \item a clean-room Tier 2 denoiser and a monocular Tier 3 reconstruction chain with frame-0 alignment.
\end{itemize}

\section{Approach}
\label{sec:approach}

\paragraph{Base recipe.}
Every policy is trained with the official FlashSAC configuration for 350\,M environment steps with 2048 parallel worlds and the official nine-stage curriculum, in which the object assist force decays to zero.
The replay buffer is 10\,M transitions on 48\,GB GPUs and 4\,M on a 24\,GB GPU.
\emph{Variant A} only changes two values: the object position reward width ($\sigma$ 0.2\,m $\rightarrow$ 0.05\,m) and the object orientation termination (1.5 $\rightarrow$ 0.6\,rad).

\paragraph{Object-orientation reward (\rtho).}
The official reward has twelve terms, and the only object term is on position.
We add a thirteenth term
\begin{equation}
\label{eq:ori}
r_{\mathrm{ori}} = w \exp\!\left(-\theta_t^2 / \sigma^2\right),
\end{equation}
where $\theta_t$ is the geodesic angle between the simulated and the reference object orientation, with $w=1.0$ and $\sigma=0.4$\,rad in all curriculum stages.
The term is set through environment variables, so checkpoints keep the official configuration and the official evaluation code reads them unchanged.

\paragraph{Selection.}
Each candidate is evaluated in 4096 deterministic worlds against the reference of its own tier; Tier 2 and Tier 3 never use Tier 1 trajectories.
A candidate replaces the policy already in the list for that episode only if ADD-AUC rises by at least 0.03, neither success rate drops by more than 0.02, and MPPE rises by at most 0.5\,cm (Tier 1 uses analogous rules on SR and MPPE).
A new list is uploaded only if the mean ADD-AUC rises by at least 0.010 on the exported rollout and 0.005 on 4096 worlds, and if the predicted board score passes the next team by 0.002.
The prediction multiplies the smaller of the two gains by 0.21, the ratio of board to own-reference gain we measured on a Tier 2 upload (+0.0039 against +0.0187).

\paragraph{Tier 2 denoiser.}
The noisy trajectories are cleaned before retargeting by a hand-designed filter whose parameters come from Tier 2 data only.
It holds static segments on a support surface, bridges motion segments with Brownian-bridge interpolation, fixes flips and height, and removes penetration by the smallest rotation before smoothing.

\paragraph{Tier 3 reconstruction.}
From the video we segment the person and the object with Grounding DINO and SAM\,2~\cite{liu2024groundingdino,ravi2024sam2}, estimate depth with MoGe~\cite{wang2025moge}, track the object with FoundationPose~\cite{wen2024foundationpose} inside a human-object reconstruction pipeline (CARI4D), and fit the SOMA body.
The camera focal length comes from a ceiling-grid calibration (866\,px instead of the default 830\,px).
Static objects slide along their viewing ray onto the support surface, a dynamics pass removes implausible object motion, and per-frame object poses are chosen among FoundationPose candidates for temporal consistency.
Finally, as the organizers allowed on 21 October, one rigid transform per episode aligns frame 0 with the Tier 1 frame-0 pose; no later Tier 1 frame is used.

\section{Experiments and Results}
\label{sec:experiments}

We trained about 370 full-length policies on one RTX\,4090, one L40S and one L40; a policy takes about two hours of training and ten minutes of evaluation and export.
Table~\ref{tab:board} lists the public board scores.
On Tier 3 the frame-0 alignment and the reconstruction updates raised the board ADD-AUC from 0.106 to 0.282 and SP-SR from 0.267 to 0.467 between 18 and 29 October; on Tier 1 the per-episode swaps lowered MPPE from 24.5\,cm (13 October) to 8.47\,cm.

\begin{table}[t]
  \centering
  \small
  \begin{tabular}{lcccc}
    \toprule
    Tier & ADD-AUC & SP-SR & MP-SR & MPPE (cm) \\
    \midrule
    1 & 0.635 (11) & 0.867 (9) & 0.700 (11) & 8.47 (9) \\
    2 & 0.518 (6)  & 0.733 (5) & 0.533 (6)  & 15.32 (9) \\
    3 & 0.282 (5)  & \textbf{0.467 (1)} & 0.100 (3) & 36.60 (4) \\
    \bottomrule
  \end{tabular}
  \caption{Public board scores and ranks of maniphys.ai on 3 November 2026 (Kaggle tie rule).}
  \label{tab:board}
\end{table}

\begin{table}[t]
  \centering
  \small
  \begin{tabular}{llccc}
    \toprule
    Tier & Ep. & Variant A & +\,\rtho & Swap \\
    \midrule
    2 & 0  & 0.400 & 0.774 & yes \\
    2 & 10 & 0.397 & 0.594 & yes \\
    2 & 18 & 0.521 & 0.855 & yes \\
    2 & 21 & 0.516 & 0.503 & no \\
    2 & 26 & 0.465 & 0.715 & yes \\
    3 & 0  & 0.528 & 0.795 & yes \\
    3 & 12 & 0.507 & 0.672 & no (MP-SR) \\
    3 & 18 & 0.626 & 0.826 & yes \\
    3 & 24 & 0.731 & 0.835 & yes \\
    3 & 26 & 0.579 & 0.812 & yes \\
    \bottomrule
  \end{tabular}
  \caption{4096-world ADD-AUC against the tier's own reference, with and without the orientation term (same seed, buffer and reference).}
  \label{tab:r3o}
\end{table}

\paragraph{Orientation term.}
Table~\ref{tab:r3o} shows part of the comparison.
Over the 12 tracked Tier 2 and Tier 3 episodes we tested, ADD-AUC rose on 11, by 0.08 to 0.37; 10 met all swap conditions.
It also helped on Tier 1 (ep24 0.580 $\rightarrow$ 0.948).
It did not help episodes that lose the object, and it is not free of risk: on Tier 2 ep21 the 4\,M run lost MP-SR (0.965 $\rightarrow$ 0.316), and the 10\,M run of Tier 3 ep26 collapsed to zero success, so that episode keeps its 4\,M run.

\paragraph{Per-episode fixes on Tier 1.}
On episodes that never track we tested one change at a time. A larger grasp-closure weight (5 $\rightarrow$ 20) lowered the 4096-world MPPE of ep9 from 40.5 to 17.0\,cm, an earlier reset when the palm leaves the reference (0.15 $\rightarrow$ 0.08\,m) lowered ep25 from 114.1 to 68.7\,cm, and the orientation term on the base recipe raised SP-SR of ep17 from 0 to 0.932. None of them helped the other lost episodes.

\paragraph{Own reference versus board.}
Gains against the own reference transferred only partly to the boards: a Tier 2 list with +0.0187 own-reference ADD-AUC gained +0.0039 on the board, and an earlier list lost 0.0029.
On Tier 3, 12 multi-phase successes against our reconstruction became 3 on the board.

\balance

\section{Conclusion}
\label{sec:conclusion}

One reward term on object orientation and strict per-episode selection gave most of our gains, with no change to the simulator, network or schedule.
The main remaining error is on episodes that lose the object at first contact and, on Tier 3, the distance between the reconstruction and the true trajectory.
At about two GPU-hours per policy on commodity GPUs, every finished run becomes a reusable candidate, so new episodes need no retuning of the rest.

%%%%%%%%%%%%%%%%%%%%%%%%%%%%%%%%%%%%%%%%%%%%%%%%%%%%%%%%%%%%%%%%%%%%%%%%

%%% References, acknowledgments and appendices do not count toward
%%% the two-page limit.
\bibliographystyle{plainnat}
\bibliography{references}

%%%%%%%%%%%%%%%%%%%%%%%%%%%%%%%%%%%%%%%%%%%%%%%%%%%%%%%%%%%%%%%%%%%%%%%%

%%% Optional appendix: starts on a new page and does not count toward
%%% the two-page limit. Delete this part if you do not need it.
\end{document}
